\documentclass[12pt,a4paper,oneside]{ctexart}
\usepackage{amsmath,amsthm,amssymb,bm,graphicx,hyperref,mathrsfs}

\title{\textbf{例题}}
% \author{黄伟宁}
\date{\today}
\linespread{1.5}
\newcounter{problemname}
\newenvironment{problem}{\stepcounter{problemname}\par\noindent\textbf{题目\arabic{problemname}. }}{\\\par}
\newenvironment{solution}{\par\noindent\textbf{解答. }}{\\\par}
\newenvironment{note}{\par\noindent\textbf{题目\arabic{problemname}的注记. }}{\\\par}
\usepackage{geometry}
\geometry{left=2.54cm,right=2.54cm,top=3.18cm,bottom=3.18cm}

\begin{document}

\maketitle
% \begin{problem}
% 这里是题目.
% \end{problem}

% \begin{solution}
%     这里是解答.
% \end{solution}

% \begin{note}
%     这里是注记.
% \end{note}

% \begin{problem}
% \end{problem}
\begin{problem}
如图,在$ \triangle OAB $ 中, $ \overrightarrow{OC}=\frac{1}{4}\overrightarrow{OA},\overrightarrow{OD}=\frac{1}{2}\overrightarrow{OB} $,$ AD $与$ BC $交于点$ M $. 设$ \overrightarrow{OA}=\boldsymbol{a},\overrightarrow{OB}=\boldsymbol{b} $.

(1)用$ \boldsymbol{a},\boldsymbol{b} $表示$ \overrightarrow{OM} $;

(2)已知在线段$ AC $上取一点$ E $, 在线段$ BD $上取一点$ F $, 使$ EF $过点$ M $. 设$ \overrightarrow{OE}=p\overrightarrow{OA},\overrightarrow{OF}=q\overrightarrow{OB} $,证明:$ \frac{1}{7p}+\frac{3}{7q}=1 $.
\end{problem}

\begin{problem}
    如图, 设点$ O $在$ \triangle ABC $内部, 且有$ \overrightarrow{OA}+2\overrightarrow{OB}+3\overrightarrow{OC}=\boldsymbol{0} $,求$ \triangle ABC $的面积与$ \triangle AOC $的面积的比值. 
\end{problem}

\begin{problem}
    已知$ O $为$ \triangle ABC $的内心, $ \cos \angle BAC=\frac{1}{3} $,且满足$ \overrightarrow{AO}=x\overrightarrow{AB}+y\overrightarrow{AC} $,则$ x+y $的最大值为\underline{\qquad\qquad}.
\end{problem}

\begin{problem}
    设$ O $为$ \triangle ABC $的外心, 若$ \overrightarrow{AO}=\overrightarrow{AB}+2\overrightarrow{AC} $,则$ \sin \angle BAC $的值为\underline{\qquad\qquad}. 
\end{problem}

\begin{problem}
    在$ \triangle ABC $中, $ O $为外心, $ G $为重心, $ H $为垂心. 证明:$ O,G,H $三点共线, 且$ |OG|=\frac{1}{2} |GH| $.
\end{problem}

\begin{problem}
    如图,设$ A_1A_2A_3A_4 $为$ \odot O $的内接四边形, $ H_1,H_2,H_3,H_4 $依次为$ \triangle A_2A_3A_4,\triangle A_3A_4A_1,$ $\triangle A_4A_1A_2,\triangle A_1A_2A_3 $的垂心. 证明:点$ H_1,H_2,H_3,H_4 $在同一圆上, 并定出该圆的圆心. 
\end{problem}

\begin{problem}
    在$ \triangle ABC $, 设角$ A,B,C $的对边分别为$ a,b,c $, 证明:

    (1)射影定理:$ a=b\cos C+c\cos B $,$ b=c\cos A+a\cos C $,$ c=a\cos B+b\cos A $;

    (2)余弦定理:$ a^2=b^2+c^2-2bc\cos A $,$ b^2=a^2+c^2-2ac\cos B $,$ c^2=a^2+b^2-2ab\cos C $.
\end{problem}

\begin{problem}
    已知$ |\boldsymbol{a}|=4,|\boldsymbol{b}|=2,\langle\boldsymbol{a},\boldsymbol{b}\rangle=\frac{\pi}{3},\boldsymbol{c}^2-\boldsymbol{a}\cdot\boldsymbol{c}=5 $, 求$ |\boldsymbol{b}-\boldsymbol{c}| $的最小值. 
\end{problem}

\begin{problem}
    设圆$ O $的半径为3, 其一条弦$ AB=4 $, $ P $为圆上任意一点, 则$ \overrightarrow{AB}\cdot\overrightarrow{BP} $的最大值为 (\qquad)

    A. 0\qquad B. 1\qquad C. 3\qquad D. 4\qquad
\end{problem}

\begin{problem}
    如图, 在$ \text{Rt}\triangle ABC $中, 已知$ BC=a $. 若长为$ 2a $的线段$ PQ $以$ A $为中点, 问:当$ \overrightarrow{PQ} $与$ \overrightarrow{BC} $的夹角$ \theta $取何值时, $ \overrightarrow{BP}\cdot\overrightarrow{CQ} $的值最大？试求出这个最大值. 
\end{problem}

\end{document}